% !TEX encoding = UTF-8 Unicode
% !TeX program = xelatex
% !TEX TS-program = XeLaTeX

\documentclass[12pt, a4paper]{article}
%\documentclass[12pt, a5paper]{article}
\usepackage{mathtools}
\usepackage{unicode-math}
\usepackage{comment}
\usepackage{wrapfig}
\usepackage{xfrac}
\usepackage{tikz}
\usetikzlibrary{arrows, decorations.markings}

\usepackage[pass]{geometry}
%\geometry{legalpaper, landscape, margin=1in}
\usepackage{fullpage}
\usepackage{setspace}
\usepackage[localise]{xepersian}
\setstretch{1.195}
\settextfont[Scale=1]{XB Niloofar}
\setlatintextfont[Scale=1]{XB Niloofar}%{Junicode} %Junicode
\setdigitfont{XB Niloofar}%{XB Zar}

\setlength{\parindent}{0ex}

\definecolor{answercolor}{HTML}{0080ff}
\newcommand{\vsp}{\vspace{34pt}}
\newcommand{\پاسخ}[2]{#1}
%\thispagestyle{empty}
%\newcommand{\پاسخ}[2]{\color{answercolor}\textbf{پاسخ: }#2\color{black}\vspace{-0.3cm}}

\begin{document}
		\tikzset{
		treenode/.style = {align=center, inner sep=0pt, text centered,
			font=\sffamily},
		arn_n/.style = {treenode, circle, white, font=\sffamily\bfseries, draw=black,
			fill=black, text width=1.5em},% arbre rouge noir, noeud noir
		arn_r/.style = {treenode, circle, red, draw=red, 
			text width=1.5em, very thick},% arbre rouge noir, noeud rouge
		arn_s/.style = {treenode, circle, black, draw=black, 
			text width=1.5em, very thick},% arbre rouge noir, noeud rouge
		arn_x/.style = {treenode, rectangle, draw=black,
			minimum width=0.5em, minimum height=0.5em, very thick}% arbre rouge noir, nil
	}
	\begin{tikzpicture}[remember picture,overlay]	\node[anchor=north,yshift=-1.5cm] at (current page.north) {\rl{به نام خدا}};
	\end{tikzpicture}%

\vspace{-1cm}
{\centering
	\textbf{کوئیز پنجم درس ساختمان داده‌ها و الگوریتم‌ها}\\
مدرس: مهدی صفرنژاد       - ۳۰ اردیبهشت ۹۸\\
\raggedleft\small
نام و نام خانوادگی: \\
}
\vspace{-0.5cm}
\پرش‌متوسط\پرخط‌ا\\
در سؤالاتی که توضیح خواسته شده، یک پاسخ کوتاه (حدوداً بین یک تا دو خط) مدنظر است. در سؤالاتی که خواستۀ مسئله توضیحی نیست، جواب نهایی کفایت می‌کند و نوشتن مراحل رسیدن به جواب نیاز نیست.
\پاراگراف{سؤال ۱.}\sloppy
در الگوریتم تقسیم با روش نیوتون از 
$\Theta(d)$
ضرب استفاده می‌شود، ولی اگر زمان اجرای الگوریتم اجرا شده برای ضرب 
$T(d)$
باشد، زمان اجرای تقسیم به جای 
$\Theta(d.T(d))$
برابر با
$\Theta(T(d))$
هست. به اختصار بگویید چطور چنین چیزی ممکن است؟ $d$ تعداد ارقام دو عدد است.

\پاسخ{\vsp\vsp}{
با توجه به خطای هر مرحله از روش نیوتون می‌توان ضربِ هر مرحله را فقط تا دقت لازم انجام داد. در این صورت ضرب‌های اول هزینۀ کمتری دارند که به تدریج و به صورت یک سری هندسی افزایش پیدا می‌کند. در این صورت مجموع هزینه‌ها، ضریب ثابتی از هزینۀ آخرین ضرب خواهد شد.

}

\پاراگراف{سؤال ۲.}
درخت پسوندی \متن‌تایپ{adad\$} را رسم کنید.

\پاسخ{\vsp\vsp}{
	\begin{center}
	\begin{tikzpicture}[-,>=stealth',level/.style={sibling distance = 3cm/#1,
		level distance = 1.5cm}] 
	\node [arn_s] { } 
	child{ node [arn_s] {} 
	edge from parent node[above left] {\texttt{\$}} }
	child { node [arn_s] {}
		child{ node [arn_s] {} edge from parent node[left] {\texttt{\$}} 
}
child{ node [arn_s] {} edge from parent node[right] {\texttt{ad\$}} 
}
	edge from parent node[left] {\texttt{ad}} }
	child { node [arn_s] {}
		child{ node [arn_s] {} edge from parent node[left] {\texttt{\$}} 
		}
		child{ node [arn_s] {} edge from parent node[right] {\texttt{ad\$}} 
		}
	edge from parent node[above right] {\texttt{d}} }

	; 
	\end{tikzpicture}
\end{center}
	
}

\پاراگراف{سؤال ۳.}
آرایه پسوندی و آرایۀ کمکی طولانی‌ترین پیشوند مشترک (LCP) رشتۀ \متن‌تایپ{adad\$} را بنویسید.

\پاسخ{\vsp\vsp}{
پاسخ در جدول زیر آمده است. نوشتن ستون پسوندها برای کسب نمره نیاز نیست.
	\begin{center}
		\begin{tabular}{|l|c|c|}
			\hline
			پسوندها & آرایه پسوندی & LCP \\
						\hline
									\hline
			\texttt{\$}      & 4            & 0   \\
						\hline
			\texttt{ad\$}    & 2            & 2   \\
						\hline
			\texttt{adad\$}  & 0            & 0   \\
						\hline
			\texttt{d\$}     & 3            & 1   \\
						\hline
			\texttt{dad\$}   & 1            & 0  \\
						\hline
		\end{tabular}
	\end{center}

}
	
\پاراگراف{سؤال ۴.}
تبدیل باروز ویلر \متن‌تایپ{adad\$} چیست؟

\پاسخ{\vsp\vsp}{
	 \متن‌تایپ{dd\$aa}
	
}

\پاراگراف{سؤال ۵.}
عبارت
$a^{b+c}\cdot c^{(b-a)\cdot(d^e)}$
را به صورت پسوندی بنویسید.

\پاسخ{\vsp\vsp}{
\lr{\texttt{a b c + \^{} c b a - d e \^{} * \^{} *}}
	
}

\پاراگراف{سؤال ۶.}
درخت دکارتی لیست
$[10, 3, 5, 7, 4]$
را رسم کنید.

\پاسخ{\vsp\vsp}{
	\begin{center}
	\begin{tikzpicture}[-,>=stealth',level/.style={sibling distance = 2cm/#1,
		level distance = 1cm}] 
	\node [arn_s] at (0, 0) {3}
	child{ node [arn_s] at (0, 0) {10} }
	child { node [arn_s] at (4, 0) {4}
		child{ node [arn_s] at (-4, 0) {5} 
			child{ node [arn_s] at (2, 0) {7} 
			}
		}
	 }	
	; 
	\end{tikzpicture}
\end{center}
}

\پاراگراف{سؤال ۷.}
برای حل مسئلۀ پرسش کمینۀ یک محدوده ایدۀ توان‌های ۲ و درخت پاره‌خطی وجود دارد. با توجه به هر کدام از معیارهای حافظه، زمان پاسخ به پرسش (query) و قابلیت تغییر یکی از اعضای آرایه کدام از این دو راه‌حل بهتر هستند؟

\پاسخ{\vsp\vsp}{
	از لحاظ حافظه، درخت پاره‌خطی بهتر است زیرا حافظه آن $O(n)$ و حافظۀ ایدۀ توان‌های ۲ $O(n\log n)$ است. از لحاظ زمانِ پاسخ به پرسش، ایدۀ توان‌های ۲ بهتر است زیرا به هر پرسش در $O(1)$ پاسخ می‌دهد و درخت پاره‌خطی در $O(\log n)$. از لحاظ قابلیت تغییر یکی از اعضای آرایه، درخت پاره‌خطی این امکان را دارد ولی ایده توان‌های ۲ این قابلیت را ندارد. برای نمره‌دهی، نوشتن دلایل نیاز نیست و ذکر داده‌ساختار مناسب‌تر از هر دید کفایت می‌کند.
	
}

\پاراگراف{سؤال ۸.}
آیا تابع $f$ یک داده‌ساختار افزوده (مثل درخت AVL افزوده یا درخت پاره‌خطی) می‌تواند «میانه عناصر» باشد؟ چرا؟

\پاسخ{\vsp\vsp}{
	خیر. زیرا مقدار $f$ یک رأس را نمی‌توان از مقادیر $f$ دو فرزند آن به دست آورد.
	
}

\پاراگراف{سؤال ۹.}
در داده‌ساختار مجموعه‌های مجزا، rank چه چیزی را نشان می‌دهد؟

\پاسخ{\vsp\vsp}{
	rank
	همان ارتفاع است در حالتی که فشرده‌سازی مسیر نداشته باشیم. به عبارت دیگر rank حداکثر ارتفاع ممکن زیردرخت است.
	
}

\پاراگراف{سؤال ۱۰.}
یک داده‌ساختار پشته داریم که علاوه بر \متن‌لاتین{pop()} (با هزینه $c$) و \متن‌لاتین{push(x)} (با هزینه $c$) دارای عمل \متن‌لاتین{multipop(k)} (با هزینه $kc$) باشد که معادل $k$ بار اجرای \متن‌لاتین{pop()} است. می‌خواهیم با تحلیل تابع پتانسیل ثابت کنیم هزینه سرشکن هر عمل $O(1)$ است. یک تابع پتانسیل مناسب برای این کار پیشنهاد دهید.

\پاسخ{\vsp\vsp}{
تعداد اعضای پشته ضرب در $c$. در صورتی که ضرب در $c$ لحاظ نشود، نیمی از نمره داده می‌شود.
	
}
\vspace{1pt}
\begin{flushleft}
موفق باشید
\end{flushleft}
\end{document}